% Generated draft: requires mathematical and source review.
\documentclass[11pt]{article}
\usepackage[margin=25mm]{geometry}
\usepackage{fontspec}
\setmainfont{Latin Modern Roman}
\usepackage{amsmath,amssymb,mathtools,unicode-math}
\setmathfont{Latin Modern Math}
\usepackage{enumitem,needspace,xurl,hyperref}
\setlength{\parindent}{0pt}
\setlength{\parskip}{5pt}
\setlength{\emergencystretch}{3em}
\newcommand{\N}{\mathbb{N}}
\newcommand{\Z}{\mathbb{Z}}
\newcommand{\Q}{\mathbb{Q}}
\newcommand{\R}{\mathbb{R}}
\newcommand{\C}{\mathbb{C}}
\newcommand{\F}{\mathbb{F}}
\newcommand{\A}{\mathbb{A}}
\newcommand{\PP}{\mathbb{P}}
\newcommand{\E}{\mathbb{E}}
\newcommand{\eps}{\varepsilon}
\newcommand{\dd}{\,\mathrm{d}}
\begin{document}
\section*{UnsolvedMath ID: 30002439}
\textbf{Source problem code:} catalogue entry
\subsection{Definitions and mathematical statement}
The following formulation is conditional on interpreting the morphism as defined over \(\Q\) and \(H\) as the standard multiplicative projective height; neither convention is explicitly specified in the supplied statement.

Let \(n,m\) be projective dimensions and let \(d\geq 1\) be an integer. Let
\[
f:\PP^n_{\Q}\longrightarrow\PP^m_{\Q},\qquad
[x_0:\cdots:x_n]\longmapsto[F_0(x):\cdots:F_m(x)],
\]
where the \(F_i\in\Q[X_0,\ldots,X_n]\) are homogeneous forms of degree \(d\) with no common zero in \(\PP^n_{\overline{\Q}}\). Here degree \(d\) refers to the degree of the defining forms. Assume that \(f\) is a closed immersion, meaning that it identifies \(\PP^n_{\Q}\) isomorphically with a closed subscheme of \(\PP^m_{\Q}\).

For \(y\in\PP^m(\Q)\), choose primitive integral homogeneous coordinates \((y_0,\ldots,y_m)\), and define
\[
H(y)=\max_{0\leq i\leq m}|y_i|.
\]
For real \(B\geq1\), put
\[
n_f(B)=\#\{x\in\PP^n(\Q):H(f(x))\leq B\}.
\]
The supplied question asks whether, for every \(\varepsilon>0\),
\[
n_f(B)=O_{d,n,\varepsilon}\!\left(B^{(n+1)/d+\varepsilon}\right)
\qquad(B\longrightarrow\infty)
\]
holds for every such closed immersion.

The displayed dependence suggests that the implied constant may depend only on \(d,n,\varepsilon\), independently of \(f\), its coefficients, and \(m\). The source does not explicitly specify this uniformity or whether the threshold beyond which the estimate holds must also be uniform. Consequently, a fully quantified uniform formulation cannot be fixed from the supplied record alone. The base field, height convention, and intended ranges of \(n,m\), including zero-dimensional boundary cases, also require confirmation. These unresolved qualifications are part of this statement-only draft.
\subsection{Short English statement}
The question asks whether a closed immersion of projective space defined by degree-\(d\) forms has at most \(O_{d,n,\varepsilon}(B^{(n+1)/d+\varepsilon})\) rational preimages whose images have height at most \(B\), for every \(\varepsilon>0\). The intended height, base field, dimension ranges, and precise uniformity remain to be confirmed. The supplied registry labels the problem open; this draft does not independently verify that status.
\subsection{Sources}
\begin{itemize}
\item Supplied local UnsolvedMath registry record, permanent ID \(30002439\), title \emph{Height Counts for Closed Projective Immersions}, in the Oberwolfach Reports Open Problems collection. Exact source problem code: \texttt{catalogue entry}. The record gives proposed year \(2013\) and names no proposer.
\item Supplied source citation: \emph{Analytic Number Theory} (2014), \emph{Oberwolfach Reports}, pp.\ 2963--3037, DOI \texttt{10.4171/owr/2013/51}. Original OWR report: \url{https://doi.org/10.4171/owr/2013/51}. The supplied background describes these reports as documenting mathematical workshops and their open-problem discussions. The citation year \(2014\) and the proposed year \(2013\) are retained as supplied.
\item Supplied literature source, labeled ``Salberger rational point bounds'': \url{https://doi.org/10.1016/j.ansens.2004.10.005}. No more specific bibliographic information or theorem statement is supplied.
\item Literature assessment dated exactly \texttt{2026-08-21}: the supplied curation labels the problem open, says that the sharper closed-immersion height exponent is still presented as a conjecture, and reports that no general proof was found. This is supplied dated literature triage derived from the OWR curation, not an independently verified assessment. The record lists the original OWR report and the Salberger reference above as sources checked. Time-sensitive status and cited claims require independent verification.
\item The registry assigns the statement status \texttt{corrected\_verified} and says that notation, typography, and wording were normalized without changing mathematical content. This is a supplied editorial designation, not mathematical verification performed for this draft.
\end{itemize}
\end{document}
